% ############################################################################
% ##  SUPERSEDED 2026-08-03 -- NOT COMPILED BY ANY DOCUMENT. DO NOT EDIT.
% ##
% ##  This file is a per-theorem rendering from an ABANDONED architecture in
% ##  which paper/main.tex would \input one file per result.  paper/main.tex
% ##  DOES NOT EXIST.  The manuscript is paper/qmd/{main,supplement}.qmd, and
% ##  both \input the theory corpus directly.  Verified 2026-08-03: neither
% ##  .qmd references sections/theorems.
% ##
% ##  The text below is ROUND-2 vintage against a corpus now at dispatch 25.
% ##  66-78% of its substantive lines no longer match the corpus -- that is old
% ##  text the corpus has since revised, NOT unique work.  Every \label here is
% ##  also declared in the live corpus.
% ##
% ##  SOURCE OF RECORD:
% ##    statements  quality_reports/theory/panel-npiv-discrete/results.tex
% ##    assumptions quality_reports/theory/panel-npiv-discrete/assumptions.tex
% ##    proofs      quality_reports/theory/panel-npiv-discrete/proofs.tex
% ##
% ##  Kept rather than deleted because nothing in this repo is committed.
% ##  Full rationale: paper/sections/theorems/_SUPERSEDED.md
% ############################################################################

% ============================================================================
%  prp:nophifree / rem:nophifree / cor:udep -- the impossibility result, its
%  interpretation, and the corollary that forces the debiasing weight to vary
%  with the residual.
%  Project: panel-npiv-discrete.  Target of plan Sec.2.7; memo Sec.10 Step 3.
%
%  Source of record: quality_reports/theory/panel-npiv-discrete/{results,proofs}.tex
%  Assumptions cited by tag; see .../assumptions.tex.
%  Requires \IR, \un, \upmodels, \vdagger from preamble.tex.
%
%  DEPENDENCIES that must appear EARLIER in the compiled document:
%    def:model     -- eq:trialres, the residual at trial values;
%    def:spaces    -- V, nu_g (eq:nudef), L^2(pi) = V (+) constants;
%    lem:nusigma   -- clauses (iii) and (iv), from lem-adjoint.tex.
%  Conventions (P1), (P5) are stated once in proofs.tex.
%  Assumption tags used: prp:nophifree -- C1(ii) ONLY;
%                        cor:udep      -- A1, A1b(i), A1b(iv), N1, C1(ii).
%  ROUND-8: A1b(i) STRUCK from prp:nophifree, statement and proof.  Step 3
%  upgrades "orthogonal to V" to "zero" using def:spaces (the decomposition
%  L^2(pi) = V (+) constants) together with int nu_g = 0 from lem:nusigma(iv);
%  varphi_dagger occurs nowhere in the proposition.  A1b(i) is KEPT on cor:udep,
%  whose clause (a) is "first-order insensitive AT varphi_dagger" -- a genuine
%  well-posedness use of the base point.
%
%  ENVIRONMENT NOTE.  The formal statement is a PROPOSITION with a proof;
%  rem:nophifree is the remark that follows it, per house style.  Plan Sec.2.7
%  registers this target under the label rem:nophifree, which is retained.
%
%  THIS FILE IS AN ALTERNATIVE RENDERING of prp:nophifree, cor:udep and
%  rem:nophifree, carrying the SAME \label keys as
%  quality_reports/theory/panel-npiv-discrete/results.tex.  The two must never
%  be \input into the same document.
% ============================================================================

\begin{proposition}\label{prp:nophifree}
A moment is free of the structural function exactly when the two periods' $g$-tilted marginals of the
endogenous regressor coincide.
Suppose Assumption C1(ii) holds and let $g$ be a random variable satisfying
\begin{equation}\label{eq:gproviso}
E|g|<\infty,
\qquad
E\big|g\,h(X_j)\big|<\infty\ \ (j=1,2)\ \text{for every }h\in\mathcal V,
\qquad
E\big|g\,\Delta U(\theta,0)\big|<\infty .
\end{equation}
Then the map $\varphi\mapsto E[g\,\Delta U(\theta,\varphi)]$ is constant on $\mathcal V$, for one and
hence every $\theta\in\IR^{p}$, if and only if
\begin{equation}\label{eq:tilted}
E[g\mid X_2=x]\,f_{X_2}(x)\;=\;E[g\mid X_1=x]\,f_{X_1}(x)
\qquad\text{for almost every }x\in\mathcal X,
\end{equation}
that is, if and only if $\nu_g\equiv0$.
\end{proposition}

\begin{proof}
The proof differences out the constant, identifies the increment with a linear functional, and uses the
decomposition $L^{2}(\pi)=\mathcal V\oplus\{\text{constants}\}$ together with
$\int_{\mathcal X}\nu_g\,dx=0$ from \cref{lem:nusigma}(iv) to upgrade ``orthogonal to $\mathcal V$'' to
``zero''. \emph{It works on increments
rather than on the three separate terms of $\Delta U(\theta,\varphi)$, because $\Delta Y$ and $\Delta Z$
are not functions of $X_1$ or $X_2$ and so \eqref{eq:gproviso} does not make $E[g\Delta Y]$ or
$E[g\Delta Z']$ separately finite.}

\emph{Step 0: $\nu_g$ is well defined and $\nu_g/\pi\in L^{2}(\pi)$.}
$E|g|<\infty$ makes the conditional expectations in \eqref{eq:nudef} well defined and gives
$\int|\nu_g|dx\le2E|g|<\infty$, so $\Xi:=\nu_g/\pi\in L^{1}(\pi)$ by C1(ii). For $h\in\mathcal V$,
Jensen's inequality and \eqref{eq:gproviso} give $\int|\nu_gh|dx\le\sum_jE|g\,h(X_j)|<\infty$; since
$\Xi\in L^{1}(\pi)$ and any $h\in L^{2}(\pi)$ is $h_0+\int h\,d\pi$ with $h_0\in\mathcal V$, in fact
$\Xi h\in L^{1}(\pi)$ for every $h\in L^{2}(\pi)$. \emph{This forces $\Xi\in L^{2}(\pi)$}: the
functionals $T_n(h):=\int\Xi\un\{|\Xi|\le n\}h\,d\pi$ are bounded with
$\|T_n\|=\|\Xi\un\{|\Xi|\le n\}\|_{L^{2}(\pi)}$ and converge pointwise on $L^{2}(\pi)$ by dominated
convergence with dominating function $|\Xi h|$, so Banach--Steinhaus bounds $\sup_n\|T_n\|$ and monotone
convergence gives $\|\Xi\|_{L^{2}(\pi)}<\infty$. \emph{Without this, ``$\nu_g/\pi\in\mathcal V^{\perp}$''
in Step~2 is not well formed.} Note $h\equiv1\notin\mathcal V$, which is why $E|g|<\infty$ is a separate
clause of \eqref{eq:gproviso} rather than a consequence of the second.

\emph{Step 1: the increment.}
For $\varphi,\varphi_0\in\mathcal V$, $\Delta U(\theta,\varphi)-\Delta U(\theta,\varphi_0)=-\Delta(\varphi-\varphi_0)$
by \eqref{eq:trialres}, and both products with $g$ are integrable by the third clause of
\eqref{eq:gproviso}, so
\begin{equation}\label{eq:incrfree}
E\big[g\,\Delta U(\theta,\varphi)\big]-E\big[g\,\Delta U(\theta,\varphi_0)\big]
\;=\;-\,E\big[g\,\Delta(\varphi-\varphi_0)\big]
\;=\;-\int_{\mathcal X}\nu_g\,(\varphi-\varphi_0)\,dx
\end{equation}
by \cref{lem:nusigma}(iii). No term is evaluated on its own and $\theta$ cancels.

\emph{Step 2: constancy is orthogonality.}
$\mathcal V$ is linear, so $\varphi-\varphi_0$ ranges over all of $\mathcal V$; the map is constant if
and only if $\langle\Xi,h\rangle_{L^{2}(\pi)}=0$ for every $h\in\mathcal V$, i.e.\
$\Xi\in\mathcal V^{\perp}$, meaningful by Step~0.

\emph{Step 3: the constant is zero.}
By \cref{def:spaces}, $\Xi=\mathsf c$ for some $\mathsf c\in\IR$. By \cref{lem:nusigma}(iv) with
$h\equiv1$ --- legitimate because $E|g|<\infty$ --- $\int_{\mathcal X}\nu_g\,dx=0$, and
$\int\Xi\pi\,dx=\int\nu_g\,dx$, so $\mathsf c=0$. Hence $\nu_g\equiv0$ since $\pi>0$ by C1(ii), which is
\eqref{eq:tilted}. The converse is immediate from \eqref{eq:incrfree}.
\end{proof}

\begin{corollary}[to \cref{prp:nophifree}]\label{cor:udep}
If the debiasing weight does not depend on the residual, the debiased moment is not merely first-order
insensitive to the structural function but exactly free of it.
Suppose Assumptions A1, A1b(i), A1b(iv) and C1(ii) hold, let $L$ be differentiable in $u$ at \emph{every}
$(u,z)$ with
$\partial_uL(u,z)=c(z)$ for some $c$ such that $g:=\Delta Z+\widetilde W\,c(\Delta Z)$ satisfies
\eqref{eq:gproviso} coordinatewise, and suppose $L$ satisfies N1. Then $L(u,z)=c(z)u$, the moment
function is
$\psi=g\,\Delta U(\theta,\varphi)$ with $g$ measurable with respect to $\sigma(W,\Delta Z)$, and the
following are equivalent: \emph{(a)} $\nu_{\Delta Z+\widetilde Wc}\equiv0$; \emph{(b)}
$\varphi\mapsto E[\psi]$ is constant on $\mathcal V$; \emph{(c)} $\nu_{\widetilde Wc}=-\nu$ \emph{for the
$c$ of the hypothesis} --- equivalently, membership
$-\nu\in\mathcal N_0:=\{\nu_{\widetilde W\tilde c(\Delta Z)}:\tilde c\in\mathcal C_0\}$ is realised
\emph{at that} $c$, where $\mathcal C_0$ is the set of $\tilde c$ for which
$\Delta Z+\widetilde W\tilde c(\Delta Z)$ satisfies \eqref{eq:gproviso} coordinatewise. Moreover
$\varphi\mapsto E[\psi]$ is affine on $\mathcal V$, so every term of order two or higher in its expansion
vanishes identically \emph{whether or not} \textup{(a)--(c)} hold. Contrapositively, if
$-\nu\notin\mathcal N_0$, then every $\ell_\star$ satisfying
$\nu_{\Delta Z+\widetilde W\ell_\star}\equiv0$ has $\partial_u\ell_\star\not\equiv0$.
\end{corollary}

\begin{proof}
The proof has three steps: the normalisation pins the constant of integration; the moment collapses to a
residual-free weight; and \cref{prp:nophifree} converts first-order insensitivity into exact constancy.

\emph{Step 1: N1 forces $L(u,z)=c(z)u$.}
Fix $z$ and set $\Lambda_z(u):=L(u,z)-c(z)u$, which is differentiable at every point with
$\Lambda_z'\equiv0$; by the \emph{mean value theorem} $\Lambda_z$ is constant, so $L(u,z)=c(z)u+b(z)$.
\emph{The mean value theorem is the right theorem and the fundamental theorem of calculus, cited
through round 6, is the wrong one:} under an a.e.\ reading of the derivative hypothesis the conclusion is
false without absolute continuity, the Cantor function being the witness, which is why the hypothesis is
stated pointwise.
Evaluating at the truth and conditioning on $\Delta Z$,
$E[L(\Delta U,\Delta Z)\mid\Delta Z]=c(\Delta Z)E[\Delta U\mid\Delta Z]+b(\Delta Z)=b(\Delta Z)$, the
first term vanishing by the second clause of A1. N1 sets the left-hand side to zero, so $b\equiv0$.

\emph{Step 2: the weight is $\sigma(W,\Delta Z)$-measurable.}
Substituting $L(u,z)=c(z)u$ into
$\psi=\Delta Z\Delta U(\theta,\varphi)+\widetilde WL(\Delta U(\theta,\varphi),\Delta Z)$ gives
$\psi=g\,\Delta U(\theta,\varphi)$ with $g=\Delta Z+\widetilde Wc(\Delta Z)$, a function of
$(W,\Delta Z)$ alone since $\widetilde W=W-\varrho(\Delta Z)$. That $g$ satisfies \eqref{eq:gproviso} is
the hypothesis; $|\widetilde W|\le1$ by (P1) and A1b(iv) cover the $\Delta Z$ term, but the clauses
involving $c$ must be assumed.

\emph{Step 3: the equivalences, with (c) read at the given $c$.}
By \eqref{eq:incrfree} applied to this $g$, the increment of $\varphi\mapsto E[\psi]$ between any
$\varphi,\varphi_0\in\mathcal V$ is $-\int\nu_g(\varphi-\varphi_0)dx$, so (a) is the vanishing of every
increment and (a)$\Leftrightarrow$(b) is \cref{prp:nophifree} applied to this $g$. For
(a)$\Leftrightarrow$(c), linearity of $g\mapsto\nu_g$ gives $\nu_g=\nu+\nu_{\widetilde Wc}$, so
$\nu_g\equiv0$ holds if and only if $\nu_{\widetilde Wc}=-\nu$ \emph{at the $c$ fixed in the
hypothesis}, which is what (c) asserts.

\emph{A quantifier that must not be relaxed.} Clause (c) is \emph{not} $-\nu\in\mathcal N_0$ read
existentially: $\mathcal N_0$ is indexed by all admissible $\tilde c$, so that membership asserts only
that \emph{some} $\tilde c$ realises the equality, whereas (a) and (b) concern the \emph{given} $c$ and
the $L$ built from it. The implication (a)$\Rightarrow$($-\nu\in\mathcal N_0$) holds and is all the
contrapositive uses; the reverse fails, since a different $\tilde c$ may realise membership while the
given $c$ does not, in which case the sentence about the second-order term does not apply to this $L$.

If (b) holds then $E[\psi]$ is constant in $\varphi$, so every term of order two or higher in its
expansion vanishes. Contrapositively, suppose $-\nu\notin\mathcal N_0$ and let $\ell_\star$ satisfy
$\nu_{\Delta Z+\widetilde W\ell_\star}\equiv0$. If $\partial_u\ell_\star\equiv0$ then
$\ell_\star(u,z)=\tilde c(z)$ for some admissible $\tilde c$, the hypothesis on $\ell_\star$ is (a) at
that $\tilde c$, so $\nu_{\widetilde W\tilde c}=-\nu$ and $-\nu\in\mathcal N_0$ --- a contradiction.
Hence $\partial_u\ell_\star\not\equiv0$.
\end{proof}

\begin{remark}\label{rem:nophifree}
\Cref{prp:nophifree} is the paper's negative result, and three readings of \eqref{eq:tilted} are worth
separating.

\emph{(i) The condition, in words.} Write $\tilde f^{\,g}_j(x):=E[g\mid X_j=x]f_{X_j}(x)$ for the
$g$-tilted density of the endogenous regressor in period $j$, a signed density of total mass $E[g]$.
Condition \eqref{eq:tilted} says $\tilde f^{\,g}_1=\tilde f^{\,g}_2$: \emph{the two periods' $g$-tilted
marginals coincide.} At $g\equiv1$ this is marginal stationarity of $X$, so even the unweighted moment
$E[\Delta U]$ is free of $\varphi$ only when the regressor's marginal law does not move between the
periods. At $g=\Delta Z_j$ it is $\nu\equiv0$, the design in which the debiasing weight degenerates to
zero and the moment reduces to the differenced least-squares one. The failure of \eqref{eq:tilted} is
therefore not a pathology to be assumed away: it is what makes the panel design informative, and a design
in which it holds for the weights of interest is one in which there is nothing to debias.

\emph{(ii) The same condition, reached independently in the linear case.} In the first-differenced panel
with a conditional-mean restriction, \citet{feve2014} observe that the operator
$\varphi\mapsto E[\varphi(X_2)-\varphi(X_1)\mid W]$ annihilates \emph{every} $\varphi$ when
$(X_1,X_2,W)$ has the same law as $(X_2,X_1,W)$ --- exchangeability of the two periods given the
instrument, which in particular equates the laws of $(X_1,W)$ and $(X_2,W)$. That is \eqref{eq:tilted}
imposed simultaneously for every bounded $g=g(W)$, since
$\nu_{g(W)}(x)=\sum_{j\in\{0,1\}}g(j)\{f_{X_2,W}(x,j)-f_{X_1,W}(x,j)\}$; and their Proposition~2.1 gives
the adjoint of that operator as an integral against the same difference of joint densities over the same
reference density, as recorded at \cref{rem:adjointreconcile}(iv). \emph{The difference structure is not
a device of this paper: it is inherited from first differencing and is shared with the linear
conditional-mean design.} What \cref{prp:nophifree} adds is the one-weight-at-a-time equivalence ---
their condition is the simultaneous one, and is correspondingly stronger --- and the enlargement of $g$
from functions of $W$ to functions of $(W,\Delta Z)$, which is the class the debiasing construction
searches.

\emph{(iii) What is new here.} The difference from \citet{feve2014} is not simply that their operator is
linear. It is that $\mathcal A_{\mathrm{cond}}$ is a \emph{Fr\'echet derivative of a nonlinear map}
rather than a fixed linear conditional-expectation operator, and three consequences follow that have no
counterpart there: the operator depends on the point of linearisation, with the Lipschitz modulus of
\cref{prp:frechet}(ii); the sample operator is noisy \emph{as an operator}, which is why A6 and A7 are
stated for $\hat T$ rather than for $T$; and the moment carries a second-order term in $\varphi$, which
is \cref{cor:udep}'s subject and which a linear operator equation does not produce.

\emph{(iv) Why the negative result creates the rate requirement.} \Cref{cor:udep} closes the chain. If
the debiasing weight ignored the residual, the moment would reduce to a $\sigma(W,\Delta Z)$-measurable
weight times the residual, and orthogonality would then force \eqref{eq:tilted}; in any design where
\eqref{eq:tilted} fails for such weights, the selected $\ell_\star$ must vary with the residual. That is
what puts $\partial_u\ell_\star$ into the second-order term and is why a quadratic remainder in $\varphi$
is unavoidable within this moment class rather than an artefact of the expansion. \emph{Two caveats
belong with the claim.} \Cref{cor:udep} shows $\partial_u\ell_\star\not\equiv0$, which is
\emph{necessary} for the second-order form to be nonzero but not sufficient: the conditional expectation
$q(x_1,x_2)=E[\widetilde W\partial_u\ell_\star\mid X_1=x_1,X_2=x_2]$ could still vanish by
cancellation, and excluding that is a separate hypothesis rather than a consequence. And
\cref{prp:nophifree} is an equivalence, not an impossibility theorem: whether \eqref{eq:tilted} fails is
a property of the application and should be reported rather than asserted.
\end{remark}
